\chapter{Anforderungsanalyse}
Die Anforderungsanalyse stellt einen zentralen Bestandteil der systematischen Entwicklung technischer Systeme dar.
Ziel ist die strukturierte Erfassung, Beschreibung und Bewertung aller relevanten Anforderungen an das zu entwickelnde System.
Auf dieser Grundlage erfolgt die Auslegung der Hardware- und Softwarearchitektur sowie die Ableitung konkreter Implementierungsschritte.

Im Kontext der vorliegenden Arbeit werden die Anforderungen aus der zuvor durchgeführten Marktanalyse sowie der definierten Zielsetzung abgeleitet.
Dabei werden sowohl funktionale als auch technische und nichtfunktionale Anforderungen berücksichtigt.
Die Analyse erfolgt in strukturierter Form, um eine nachvollziehbare und überprüfbare Systemauslegung zu gewährleisten.
Die Anforderungen werden in tabellarischer Form zusammengefasst und anschließend textuell präzisiert.
Dies ermöglicht eine klare Zuordnung zwischen den einzelnen Anforderungen und den entsprechenden Systemkomponenten.